% GENERATED FILE. Edit canonical lessons, then run npm run generate:books.
% canonical-source-hash: fnv1a64:fdce762ccd26e2ff
% canonical-lessons: JA-C353-tenken-suru, JA-C353-choushi, JA-C353-shoukaki, JA-C353-dansui, JA-C353-ichiou

\chapter{Checking Things Work}
\label{ch:ja-checking-things-work}
\hlchaptermodality{353}

\begin{chapteropening}
\textbf{By the end of this chapter:} \emph{I can say what to check and what to do if something stops working.}

Complete the last lesson of chapter 353: I can say what to check and what to do if something stops working.
\end{chapteropening}

\section[\texorpdfstring{\ja{てんけんする} (tenken suru)}{てんけんする (tenken suru)}]{\ja{てんけんする} (tenken suru) — to inspect; to check over}

\label{lesson:JA-C353-tenken-suru}



\begin{quote}
Before the new one: say the Japanese for to slip out of place, then the Japanese for a tool.
\end{quote}

\subsection*{You'll want to know: \ja{てんけんする}}
\textbf{\ja{てんけんする}} — \emph{tenken suru} — \textquotedblleft{}to inspect; to check over\textquotedblright{}.

\begin{quote}
\textbf{\ja{ながく} \ja{つかったら、くるまを} \ja{てんけんした} \ja{ほうが} \ja{いいです。}} — \emph{nagaku tsukattara, kuruma o tenken shita hō ga ii desu} — \textquotedblleft{}Once a car has been used for a long time, it should be checked over.\textquotedblright{}
\end{quote}

Polite \textbf{\ja{てんけんします}}, \emph{tenkenshimasu}; a \ja{する} verb. Past (polite): \textbf{\ja{ちちが} \ja{くるまを} \ja{てんけんしました。}}, \emph{chichi ga kuruma o tenken shimashita}.

\begin{cousinweb}
\textbf{\ja{てんけんする}} (\emph{tenken suru}) is built from a Sino-Japanese noun and the native verb \emph{suru}, \textquotedblleft{}to do\textquotedblright{}. The noun, \emph{tenken}, is two Chinese characters read with their on'yomi, readings borrowed from Chinese: \emph{ten}, \textquotedblleft{}spot\textquotedblright{} (Middle Chinese \emph{temX}) and \emph{ken}, \textquotedblleft{}examination\textquotedblright{} (\emph{kjemX}) (KANJIDIC2; JMdict; Middle Chinese in Baxter's transcription).
\end{cousinweb}

\begin{grammarlens}[title={What you've built}]
One more word for saying what happens if.
\end{grammarlens}

\subsection*{Your turn}
\begin{itemize}
\raggedright
  \item \emph{Say it:} \emph{tenken suru}
  \item \emph{Say it:} \emph{tenken suru}, once more
  \item \emph{Say it:} \emph{nagaku tsukattara, kuruma o tenken shita hō ga ii desu} — Once a car has been used for a long time, it should be checked over.
  \item \emph{Recall:} say the Japanese for to slip out of place, then the Japanese for a tool, then say \emph{tenken suru} again
\end{itemize}

\begin{tcolorbox}[breakable,colback=teal!4,colframe=teal!35!black,title={Before you move on}]
What does \ja{てんけんする} mean? (\textquotedblleft{}To inspect; to check over\textquotedblright{}.) Say it once more.
\end{tcolorbox}
\section[\texorpdfstring{\ja{ちょうし} (chōshi)}{ちょうし (chōshi)}]{\ja{ちょうし} (chōshi) — how things are running; condition}

\label{lesson:JA-C353-choushi}



\begin{quote}
Before the new one: say the Japanese for a tool, then the Japanese for to inspect.
\end{quote}

\subsection*{You'll want to know: \ja{ちょうし}}
\textbf{\ja{ちょうし}} — \emph{chōshi} — \textquotedblleft{}how things are running; condition\textquotedblright{}.

\begin{quote}
\textbf{\ja{ちょうしが} \ja{なおったら、しごとに} \ja{もどって} \ja{ください。}} — \emph{chōshi ga naottara, shigoto ni modotte kudasai} — \textquotedblleft{}When you are back in shape, please come back to work.\textquotedblright{}
\end{quote}

\begin{cousinweb}
\textbf{\ja{ちょうし}} (\emph{chōshi}) is Sino-Japanese (kango): two Chinese characters read with their on'yomi, readings borrowed from Chinese. They are \emph{chō}, \textquotedblleft{}tune\textquotedblright{} (Middle Chinese \emph{dewH}, the reading for a tone or melody) and \emph{shi}, \textquotedblleft{}child\textquotedblright{} (\emph{tsiX}) (KANJIDIC2; Middle Chinese in Baxter's transcription).
\end{cousinweb}

\begin{grammarlens}[title={What you've built}]
One more word for saying what happens if.
\end{grammarlens}

\subsection*{Your turn}
\begin{itemize}
\raggedright
  \item \emph{Say it:} \emph{chōshi}
  \item \emph{Say it:} \emph{chōshi}, once more
  \item \emph{Say it:} \emph{chōshi ga naottara, shigoto ni modotte kudasai} — When you are back in shape, please come back to work.
  \item \emph{Recall:} say the Japanese for a tool, then the Japanese for to inspect, then say \emph{chōshi} again
\end{itemize}

\begin{tcolorbox}[breakable,colback=teal!4,colframe=teal!35!black,title={Before you move on}]
What does \ja{ちょうし} mean? (\textquotedblleft{}How things are running; condition\textquotedblright{}.) Say it once more.
\end{tcolorbox}
\section[\texorpdfstring{\ja{しょうかき} (shōkaki)}{しょうかき (shōkaki)}]{\ja{しょうかき} (shōkaki) — a fire extinguisher}

\label{lesson:JA-C353-shoukaki}



\begin{quote}
Before the new one: say the Japanese for to inspect, then the Japanese for how things are running.
\end{quote}

\subsection*{You'll want to know: \ja{しょうかき}}
\textbf{\ja{しょうかき}} — \emph{shōkaki} — \textquotedblleft{}a fire extinguisher\textquotedblright{}.

\begin{quote}
\textbf{\ja{ひが} \ja{ちいさかったら、しょうかきを} \ja{つかって} \ja{ください。}} — \emph{hi ga chiisakattara, shōkaki o tsukatte kudasai} — \textquotedblleft{}If the fire is small, use the fire extinguisher.\textquotedblright{}
\end{quote}

\begin{cousinweb}
\textbf{\ja{しょうかき}} (\emph{shōkaki}) is Sino-Japanese (kango): three Chinese characters read with their on'yomi, readings borrowed from Chinese. They are \emph{shō}, \textquotedblleft{}extinguish\textquotedblright{} (Middle Chinese \emph{sjew}), \emph{ka}, \textquotedblleft{}fire\textquotedblright{} (\emph{xwaX}) and \emph{ki}, \textquotedblleft{}utensil\textquotedblright{} (\emph{khijH}) (KANJIDIC2; Middle Chinese in Baxter's transcription).
\end{cousinweb}

\begin{grammarlens}[title={What you've built}]
One more word for saying what happens if.
\end{grammarlens}

\subsection*{Your turn}
\begin{itemize}
\raggedright
  \item \emph{Say it:} \emph{shōkaki}
  \item \emph{Say it:} \emph{shōkaki}, once more
  \item \emph{Say it:} \emph{hi ga chiisakattara, shōkaki o tsukatte kudasai} — If the fire is small, use the fire extinguisher.
  \item \emph{Recall:} say the Japanese for to inspect, then the Japanese for how things are running, then say \emph{shōkaki} again
\end{itemize}

\begin{tcolorbox}[breakable,colback=teal!4,colframe=teal!35!black,title={Before you move on}]
What does \ja{しょうかき} mean? (\textquotedblleft{}A fire extinguisher\textquotedblright{}.) Say it once more.
\end{tcolorbox}
\section[\texorpdfstring{\ja{だんすい} (dansui)}{だんすい (dansui)}]{\ja{だんすい} (dansui) — a water outage}

\label{lesson:JA-C353-dansui}



\begin{quote}
Before the new one: say the Japanese for how things are running, then the Japanese for a fire extinguisher.
\end{quote}

\subsection*{You'll want to know: \ja{だんすい}}
\textbf{\ja{だんすい}} — \emph{dansui} — \textquotedblleft{}a water outage\textquotedblright{}.

\begin{quote}
\textbf{\ja{だんすいに} \ja{なったら、おふろの} \ja{みずを} \ja{つかいます。}} — \emph{dansui ni nattara, ofuro no mizu o tsukaimasu} — \textquotedblleft{}If the water is cut off, we use the water in the bath.\textquotedblright{}
\end{quote}

\begin{cousinweb}
\textbf{\ja{だんすい}} (\emph{dansui}) is Sino-Japanese (kango): two Chinese characters read with their on'yomi, readings borrowed from Chinese. They are \emph{dan}, \textquotedblleft{}severance\textquotedblright{} (Middle Chinese \emph{dwanX}) and \emph{sui}, \textquotedblleft{}water\textquotedblright{} (\emph{sywijX}) (KANJIDIC2; Middle Chinese in Baxter's transcription).
\end{cousinweb}

\begin{grammarlens}[title={What you've built}]
One more word for saying what happens if.
\end{grammarlens}

\subsection*{Your turn}
\begin{itemize}
\raggedright
  \item \emph{Say it:} \emph{dansui}
  \item \emph{Say it:} \emph{dansui}, once more
  \item \emph{Say it:} \emph{dansui ni nattara, ofuro no mizu o tsukaimasu} — If the water is cut off, we use the water in the bath.
  \item \emph{Recall:} say the Japanese for how things are running, then the Japanese for a fire extinguisher, then say \emph{dansui} again
\end{itemize}

\begin{tcolorbox}[breakable,colback=teal!4,colframe=teal!35!black,title={Before you move on}]
What does \ja{だんすい} mean? (\textquotedblleft{}A water outage\textquotedblright{}.) Say it once more.
\end{tcolorbox}
\section[\texorpdfstring{\ja{いちおう} (ichiō)}{いちおう (ichiō)}]{\ja{いちおう} (ichiō) — tentatively; for the time being}

\label{lesson:JA-C353-ichiou}



\begin{quote}
Before the new one: say the Japanese for a fire extinguisher, then the Japanese for a water outage.
\end{quote}

\subsection*{You'll want to know: \ja{いちおう}}
\textbf{\ja{いちおう}} — \emph{ichiō} — \textquotedblleft{}tentatively; for the time being\textquotedblright{}.

\begin{quote}
\textbf{\ja{あめが} \ja{ふったら} \ja{こまるので、いちおう} \ja{かさを} \ja{もちます。}} — \emph{ame ga futtara komaru node, ichiō kasa o mochimasu} — \textquotedblleft{}Rain would be a nuisance, so I'll take an umbrella anyway.\textquotedblright{}
\end{quote}

\begin{cousinweb}
\textbf{\ja{いちおう}} (\emph{ichiō}) is Sino-Japanese (kango): two Chinese characters read with their on'yomi, readings borrowed from Chinese. They are \emph{ichi}, \textquotedblleft{}one\textquotedblright{} (Middle Chinese \emph{'jit}) and \emph{ō}, \textquotedblleft{}answer\textquotedblright{} (\emph{'ingH}) (KANJIDIC2; Middle Chinese in Baxter's transcription).
\end{cousinweb}

\begin{grammarlens}[title={What you've built}]
One more word for saying what happens if.
\end{grammarlens}

\subsection*{Your turn}
\begin{itemize}
\raggedright
  \item \emph{Say it:} \emph{ichiō}
  \item \emph{Say it:} \emph{ichiō}, once more
  \item \emph{Say it:} \emph{ame ga futtara komaru node, ichiō kasa o mochimasu} — Rain would be a nuisance, so I'll take an umbrella anyway.
  \item \emph{Recall:} say the Japanese for a fire extinguisher, then the Japanese for a water outage, then say \emph{ichiō} again
\end{itemize}

\begin{tcolorbox}[breakable,colback=teal!4,colframe=teal!35!black,title={Before you move on}]
What does \ja{いちおう} mean? (\textquotedblleft{}Tentatively; for the time being\textquotedblright{}.) Say it once more.
\end{tcolorbox}
